% Artículo VII, incorporación aditiva, revisión 05.
% Arquitectura autoral: energía de memoria APP--TRIT--TPK, S01.
% Antecedente recuperado: diferencial y refinamiento de 30b.
% Formalización reunida: nota COMPOSICION_VARIACIONAL_TRIT_CONEXION.md.
% Se conserva íntegra 30b; su representante antihermítico es la
% restricción unitaria del diferencial que aquí se desarrolla.
% Esta sección evalúa la corriente de una acción y una realización
% declaradas. No selecciona retrospectivamente su campo ni su peso.

\section{Realización de la corriente en coordenadas de conexión}
\label{vii:sec:composicion-corriente-conexion}

La energía de memoria asigna a cada ruta una respuesta que conserva sus
incidencias y el orden de los transportes. Para insertarla en una acción
geométrica se necesitan sus componentes respecto de las variaciones de
conexión. El paso se obtiene diferenciando el transporte que ya realiza
la ruta: la corriente de conexión es la composición de ese diferencial
con el covector energético de la sección anterior.

Los regímenes elíptico, parabólico e hiperbólico del TRIT requieren
conservar el diferencial completo antes de restringirlo a una
representación particular. Las rotaciones unitarias constituyen un
dominio de esa construcción. La ampliación siguiente conserva también
las direcciones de dilatación y transporte que admite la realización
geométrica utilizada.

\subsection{Diferencial sobre transportes invertibles}

Se mantienen las fibras, los campos y el peso de
\eqref{vii:eq:energia-memoria-discreta}. Consideremos ahora transportes
invertibles y colecciones diferenciables con
\(\dot U_a=A_aU_a\), donde \(A_a\) pertenece al espacio tangente
de la representación elegida. Escribimos
\[
 v_a=U_af_{s(a)},\qquad r_a=f_{t(a)}-v_a,\qquad
 q=\mathsf W_mr,\qquad E_m=\langle r,\mathsf W_mr\rangle.
\]

\begin{proposition}[Representante completo del diferencial]
\label{vii:prop:diferencial-invertible}
A campo y peso fijos, se cumple
\begin{equation}
 \mathrm d_UE_m(A)=-2\operatorname{Re}\sum_a
       \langle q_a,A_av_a\rangle
 =\sum_a\operatorname{Re}\operatorname{Tr}(\mathcal G_a^*A_a),
 \qquad \mathcal G_a=-2q_av_a^*.
 \label{vii:eq:gradiente-transporte-completo}
\end{equation}
Las proyecciones antihermítica y hermítica de ese representante son
\begin{equation}
 \frac{\mathcal G_a-\mathcal G_a^*}{2}
   =v_aq_a^*-q_av_a^*=\mathcal J_a,
 \qquad
 \frac{\mathcal G_a+\mathcal G_a^*}{2}
   =-(q_av_a^*+v_aq_a^*).
 \label{vii:eq:proyecciones-gradiente-memoria}
\end{equation}
La restricción a generadores antihermíticos recupera exactamente
\eqref{vii:eq:representante-corriente-discreta}.
\end{proposition}

\begin{proof}
La identidad \(\dot r_a=-A_av_a\) y la autoadjunción del peso dan
\(\dot E_m=2\operatorname{Re}\langle q,\dot r\rangle\).
Además,
\[
 \operatorname{Tr}(\mathcal G_a^*A_a)
 =-2\operatorname{Tr}(v_aq_a^*A_a)=-2q_a^*A_av_a.
\]
Las dos proyecciones se obtienen tomando el adjunto. Para el producto
real de Hilbert--Schmidt, los subespacios hermítico y antihermítico son
ortogonales; por ello la parte hermítica tiene contracción nula con
un generador antihermítico.
\end{proof}

La fórmula permite pesos que acoplan aristas: su acción se calcula en
\(q=\mathsf W_mr\) antes de separar las contribuciones de cada
incidencia. Si varían campo o peso, se conserva la expresión completa
\eqref{vii:eq:variacion-completa-memoria}, válida también para los
transportes invertibles aquí considerados.

\subsection{Representaciones cuadráticas del TRIT}

En una realización matricial bidimensional de
\(\mathbb A_\tau\), pueden tomarse
\[
 J_{+}=\begin{pmatrix}0&1\\-1&0\end{pmatrix},\qquad
 J_{-}=\begin{pmatrix}0&1\\1&0\end{pmatrix},\qquad
 J_{0}=\begin{pmatrix}0&1\\0&0\end{pmatrix}.
\]
Por multiplicación directa,
\[
 J_+^2=-I,\qquad J_-^2=I,\qquad J_0^2=0.
\]
Sus exponenciales realizan, respectivamente, los regímenes elíptico,
hiperbólico y parabólico. El primero es ortogonal para el producto
positivo fijado. El segundo preserva la forma
\(B=\operatorname{diag}(-1,1)\), pues
\(J_-^{\mathsf T}B+BJ_-=0\). El tercero representa el álgebra
de números duales. Estas representaciones fijan ejemplos de los tipos
cuadráticos; su uso en una pantalla de otra dimensión conserva el mapa
de representación correspondiente.

Un cálculo distingue las componentes del diferencial. En una arista,
tómense \(U=I\), \(f_s=(1,1)^{\mathsf T}\),
\(f_t=2f_s\) y \(\mathsf W=I\). Entonces \(q=v=f_s\),
\(\mathcal J=0\), mientras
\begin{equation}
 \mathrm dE(J_-)=-4,\qquad \mathrm dE(J_0)=-2.
 \label{vii:eq:control-regimenes-corriente}
\end{equation}
El representante completo conserva ambas respuestas. La distinción
identifica el espacio de variaciones antes de efectuar la contracción
material.

\subsection{Composición de rutas e inversión orientada}

\begin{proposition}[Transporte contragrediente del diferencial]
\label{vii:prop:pullback-invertible}
Para \(U_\gamma=U_N\cdots U_1\), sean
\(B_k=U_N\cdots U_{k+1}\) y \(B_N=I\). Se cumple
\[
 \dot U_\gamma U_\gamma^{-1}
   =\sum_kB_kA_kB_k^{-1}.
\]
El representante \(\mathcal G_\gamma\) en la fibra final induce
en la incidencia \(k\)
\begin{equation}
 \mathcal G_k=B_k^*\mathcal G_\gamma(B_k^{-1})^*.
 \label{vii:eq:transporte-contragrediente}
\end{equation}
\end{proposition}

\begin{proof}
La regla del producto, seguida de multiplicación por
\(U_\gamma^{-1}\), demuestra la primera igualdad. La ciclicidad
de la traza da
\[
 \operatorname{Re}\operatorname{Tr}
  (\mathcal G_\gamma^*B_kA_kB_k^{-1})
 =\operatorname{Re}\operatorname{Tr}
  ((B_k^*\mathcal G_\gamma(B_k^{-1})^*)^*A_k).
\]
Para \(B_k\) unitario se recupera la conjugación de
\ref{vii:prop:composicion-respuesta}.
\end{proof}

\begin{proposition}[Inversión por congruencia energética]
\label{vii:prop:inversion-congruencia}
Para una energía local por arista, la inversión se representa por
\begin{equation}
 U_{\bar a}=U_a^{-1},\qquad
 r_{\bar a}=-U_a^{-1}r_a,\qquad
 W_{\bar a}=U_a^*W_aU_a.
 \label{vii:eq:inversion-peso-general}
\end{equation}
Conserva la energía. Si \(W_a\) permanece fijo,
\[
 A_{\bar a}=-U_a^{-1}A_aU_a,\qquad
 \dot W_{\bar a}=U_a^*(A_a^*W_a+W_aA_a)U_a.
\]
La variación completa de la energía invertida coincide con la original.
\end{proposition}

\begin{proof}
La sustitución en \(r_{\bar a}^*W_{\bar a}r_{\bar a}\)
produce \(r_a^*W_ar_a\). La derivada de la inversa es
\(\dot U_a^{-1}=-U_a^{-1}A_a\); la del peso se obtiene por
la regla del producto. La igualdad de energías vale para toda la
colección, de modo que sus derivadas coinciden, incluyendo el término
\(\langle r_{\bar a},\dot W_{\bar a}r_{\bar a}\rangle\).
\end{proof}

La congruencia conserva positividad para todo transporte invertible.
En la representación unitaria se reduce a la regla de inversión de
\ref{vii:prop:inversion-respuesta}. La suma utiliza una orientación
por arista, o la mitad de la suma sobre ambas orientaciones.

\subsection{Evaluación de las componentes de corriente}

Sea \(\theta=(\theta_b)\) una variación real de las coordenadas
de conexión de la realización utilizada. El diferencial efectivo del
transporte define los operadores
\begin{equation}
 A_a(\theta)=\dot U_aU_a^{-1}
     =\sum_bB_{ab}\theta_b.
 \label{vii:eq:diferencial-conexion-transporte}
\end{equation}
Para una composición finita, sus contribuciones se obtienen por
\eqref{vii:eq:transporte-contragrediente}. Los operadores
\(B_{ab}\) son las derivadas de la representación de conexión que
realiza la ruta.

\begin{theorem}[Corriente de conexión de la acción de memoria]
\label{vii:thm:corriente-conexion-evaluada}
Para un término material \(S_{\rm mem}=\mathfrak a_m E_m\),
con \(\mathfrak a_m\) su normalización en la recta de acción,
a campo, peso y normalización fijos se tiene
\begin{equation}
 \delta S_{\rm mem}=-\frac12\sum_b\theta_b s_b,\qquad
 s_b=4\mathfrak a_m\operatorname{Re}
        \sum_a\langle q_a,B_{ab}v_a\rangle.
 \label{vii:eq:componentes-corriente-conexion}
\end{equation}
Si el emparejamiento celular entre las bases de conexión y corriente
es una matriz invertible \(M\), definida por
\(\delta S_{\rm mem}=-\tfrac12\theta^{\mathsf T}M\sigma\),
las coordenadas de la corriente son
\begin{equation}
 \sigma=M^{-1}s.
 \label{vii:eq:coordenadas-corriente-material}
\end{equation}
\end{theorem}

\begin{proof}
Se sustituye \eqref{vii:eq:diferencial-conexion-transporte} en
\eqref{vii:eq:gradiente-transporte-completo}, se multiplica por
\(\mathfrak a_m\) y se agrupan los coeficientes de
\(\theta_b\). La convención variacional \(-1/2\) fija el
factor \(4\). Comparar ambas expresiones para toda variación
\(\theta\) da \(M\sigma=s\); la invertibilidad del
emparejamiento demuestra existencia y unicidad de esas coordenadas.
\end{proof}

La fórmula realiza explícitamente la composición mencionada en
\ref{vii:subsec:realizacion-respuesta-discreta}. Para evaluar cada
componente se transporta el campo, se calcula el residuo, se aplica el
peso y se contrae con la variación de conexión de la misma realización.
El emparejamiento \(M\) conserva orientación y volumen de las celdas.
Un emparejamiento diagonal da la correspondiente división por sus
volúmenes; otras bases conservan la matriz completa.

Si la conexión interviene también en \(f\), \(\mathsf W_m\)
o \(\mathfrak a_m\), se aplica
\eqref{vii:eq:variacion-completa-memoria} y se añade
\(E_m\delta\mathfrak a_m\). Las contribuciones de todos los
términos materiales se suman antes de componer la respuesta de Cartan.

\begin{corollary}[Compatibilidad de la corriente con el refinamiento]
\label{vii:cor:refinamiento-corriente-conexion}
Para las colecciones compatibles de
\ref{vii:thm:refinamiento-respuesta}, sea \(P_m\) el mapa fijo
que lleva variaciones de conexión del nivel \(m\) al nivel
\(m+1\). Si las acciones de memoria coinciden sobre esos campos
y transportes refinados, sus covectores satisfacen
\[
 s_m=P_m^{\mathsf T}s_{m+1},\qquad
 M_m\sigma_m=P_m^{\mathsf T}M_{m+1}\sigma_{m+1}.
\]
\end{corollary}

\begin{proof}
Diferenciar la igualdad de acciones en la dirección
\(P_m\theta\) da
\(-\tfrac12\theta^{\mathsf T}s_m
 =-\tfrac12\theta^{\mathsf T}P_m^{\mathsf T}s_{m+1}\).
La arbitrariedad de \(\theta\) y las representaciones celulares
de ambos covectores prueban las dos identidades. La composición de
estos mapas itera la igualdad a todo número finito de refinamientos.
\end{proof}

\subsection{Inserción en la ecuación de Cartan}

La sección~\ref{vii:sec:corriente-respuesta-cartan} normaliza la
tres-forma material mediante
\(\delta_\omega S_m=-\tfrac12\int\delta\omega^{IJ}\wedge
\sigma_{IJ}\). El teorema anterior proporciona sus componentes
celulares para la acción de memoria y el emparejamiento declarados.
La respuesta de Cartan utiliza esa corriente junto a la cotétrada,
el parámetro de Barbero--Immirzi y el acoplamiento gravitatorio de la
misma realización.

La dependencia material respecto de la contorsión determina cómo se
resuelve esa ecuación. Para materia afín, la corriente es independiente
de la contorsión y se aplica la inversión lineal y la eliminación
cuadrática demostradas en la sección~\ref{vii:sec:corriente-respuesta-cartan}.
Para una acción de
holonomía, el diferencial se evalúa en la conexión actual y se conserva
esa dependencia al resolver la ecuación. Por ejemplo, una arista con
\(f_s=f_t=1\), peso uno y \(U(t)=e^{it}\) tiene
\[
 E(t)=2-2\cos t,\qquad E'(t)=2\sin t.
\]
La corriente depende aquí de \(t\). La fórmula completa mantiene
esa dependencia en vez de sustituirla por su valor en una conexión
particular.

Finalmente, la densidad gravitatoria se obtiene de la variación
métrica de la acción efectiva, con su campo y su normalización.
La representación de la corriente, la contorsión que produce y la
densidad resultante son las sucesivas operaciones de esa composición.
La ley de dilución posterior transporta la amplitud así evaluada en
la sección material correspondiente.
